\subsection{Solution of the population CCA problem}
Equation \eqref{eq:gen_eig_b} has the form of a generalized eigenvalue problem
\[
A\mathbf{b} = \lambda B\mathbf{b},
\]
where
\[
A = \Sigma_{21}\Sigma_{11}^{-1}\Sigma_{12}, \quad
B = \Sigma_{22}, \quad
\lambda = \rho^2.
\]
We can consider two cases here: the simple case when we only have the covariance matrices, or, the somewhat more involved case, when we have the original data matrices at our disposal.
If $\Sigma_{22}$ is non-singular, the generalized eigenvalue problem
can be converted into a standard eigenvalue problem by multiplying
\eqref{eq:gen_eig_b} from the left by $\Sigma_{22}^{-1}$, yielding

\begin{equation} \label{eq:eig_standard_b}
    (\Sigma_{22}^{-1} \Sigma_{21} \Sigma_{11}^{-1} \Sigma_{12} - \rho^2 I)\mathbf{b} = 0
\end{equation}

Equation \eqref{eq:eig_standard_b} is a standard eigenvalue problem of the form

\[
M \mathbf{b} = \rho^2 \mathbf{b},
\]

where

\[
M = \Sigma_{22}^{-1} \Sigma_{21} \Sigma_{11}^{-1} \Sigma_{12}.
\]

Thus, the canonical correlations can be obtained by solving the
eigenvalue problem for the matrix $M$. Let

\[
M \mathbf{b}_i = \lambda_i \mathbf{b}_i ,
\]

where $\lambda_i$ are the eigenvalues and $\mathbf{b}_i$ the corresponding
eigenvectors. The canonical correlations are then given by

\[
\rho_i = \sqrt{\lambda_i}.
\]

Once the vectors $\mathbf{b}_i$ are obtained, the corresponding vectors
$\mathbf{a}_i$ can be computed from equation \eqref{eq:deriv_a} as

\[
\mathbf{a}_i =
\frac{\Sigma_{11}^{-1} \Sigma_{12} \mathbf{b}_i}{\rho_i}.
\]

The vectors $\mathbf{a}_i$ and $\mathbf{b}_i$ define the canonical
directions, and the corresponding canonical variates are

\[
U_i = \mathbf{a}_i^T \mathbf{X}^{(1)}, \qquad
V_i = \mathbf{b}_i^T \mathbf{X}^{(2)} .
\]

The number of non–zero canonical correlations is at most
$\min(p,q)$.

